%%
%% Epilepsy Research — elsarticle numbered style
%%
\documentclass[final,5p,times,twocolumn]{elsarticle}

\usepackage{amssymb}
\usepackage{amsmath}
\usepackage{booktabs}      % professional table rules
\usepackage{microtype}     % better line breaking
\usepackage{hyperref}      % clickable refs in PDF
\usepackage{lineno}        % line numbers for review

\journal{Epilepsy Research}

%% ----------------------------------------------------------------
\begin{document}
\linenumbers

\begin{frontmatter}

%% ----------------------------------------------------------------
%% Title
%% ----------------------------------------------------------------
\title{Preictal Mutual Information Dynamics Bifurcate by Seizure Onset
       Lobe: A Zero-Parameter, Four-Measure Information-Theoretic
       Analysis of the HUP iEEG Cohort}

%% ----------------------------------------------------------------
%% Authors and affiliations — fill in before submission
%% ----------------------------------------------------------------
\author[aff1]{Arif Faysal Nayem\corref{cor1}}
\ead{ariffaysal.nayem@gmail.com}
\cortext[cor1]{Corresponding author}

\affiliation[aff1]{%
  organization={},
  addressline={},
  city={},
  postcode={},
  state={},
  country={}}

%% ----------------------------------------------------------------
%% Abstract
%% Guide: 01-abstract.md — 5-sentence formula, <250 words,
%% no citations, no "novel", quantified results, past tense for
%% findings, present tense for contributions.
%% ----------------------------------------------------------------
\begin{abstract}
Mutual information (MI) has been proposed as a preictal biomarker for
focal epilepsy, yet published studies have reported opposite directions
of effect: some observe that network-wide MI decreases before seizure
onset while others observe that it increases, with no mechanistic
account of the inconsistency.
%
We tested the uniform MI collapse hypothesis --- that global mean
pairwise MI decreases preictally across all patients --- in 41 patients
from the Hospital of the University of Pennsylvania (HUP) intracranial
EEG (iEEG) dataset using four sequential, zero-parameter
information-theoretic measures: rolling Kraskov-estimated global MI,
MI-weighted dynamic network graph topology, MI versus signal variance,
and Transfer Entropy (TE) directional asymmetry.
%
The uniform collapse hypothesis was not supported at the group level
(Wilcoxon signed-rank, $p = 0.41$), and patient-level analysis
revealed that this null group result arose from two
equally prevalent, biologically opposed dynamics: frontal-onset patients
exhibited preictal MI collapse (Cohen's $d \approx +0.99$), while
mesial temporal lobe (MTL)-onset patients exhibited preictal MI
amplification ($d \approx -2.65$).
%
This lobe-consistent pattern was confirmed independently by graph
topology (16 fragmentation vs.\ 15 integration among 31 non-saturated
patients; Spearman $r = -0.60$ between MI and network density) and by
MI-variance divergence (6 of 15 frontal-onset patients showed MI
collapse while signal variance simultaneously increased; frontal-onset
interictal $r(\Phi, V) = 0.40$ vs.\ MTL-onset $r = 0.58$);
Transfer Entropy directional asymmetry yielded mean AUC = 0.509 under
imprecise MNI proxy SOZ labels (a lower bound), with MTL-targeted
patients reaching mean AUC = 0.635.
%
Preictal MI dynamics are consistent with lobe-specific ictogenic mechanisms rather than a universal process;
mixed-cohort studies that pool both phenotypes will observe a null
result --- explaining two decades of contradictory findings --- and the
MI directionality phenotype constitutes a parameter-free biomarker of
seizure network type, calibrated from the interictal baseline without
any labelled seizure data.
\end{abstract}

%% ----------------------------------------------------------------
%% Research highlights (Elsevier: 3–5 items, ≤85 chars each)
%% ----------------------------------------------------------------
\begin{highlights}
  \item No evidence for uniform preictal MI collapse at group level
        ($p = 0.41$, $N = 41$ patients).
  \item Preictal MI bifurcates by lobe: collapse in frontal-onset
        ($d \approx +0.99$), amplification in MTL-onset ($d \approx
        -2.65$).
  \item Bifurcation confirmed independently by graph topology and
        MI-variance divergence across three measures.
  \item 6 of 15 frontal-onset patients show an invisible precursor:
        MI collapses while signal variance increases simultaneously.
  \item Mixed-cohort pooling mechanically produces a null result,
        resolving two decades of contradictory literature.
\end{highlights}

%% ----------------------------------------------------------------
%% Keywords
%% ----------------------------------------------------------------
\begin{keyword}
seizure prediction \sep
mutual information \sep
intracranial EEG \sep
information bottleneck \sep
network bifurcation \sep
graph fragmentation \sep
Transfer Entropy \sep
zero-parameter biomarker \sep
preictal dynamics
\end{keyword}

\end{frontmatter}

%% ================================================================
%% SECTION 1 — INTRODUCTION
%% Guide: 02-introduction.md — 3 paragraphs:
%%   P1: problem (3–4 named papers, specific gap, scope)
%%   P2: literature (5–6 papers, thematic groups, bridging sentence)
%%   P3: contribution (what remains, mechanism, outcomes, bullets,
%%       roadmap)
%% ================================================================
\section{Introduction}
\label{sec:intro}

Drug-resistant focal epilepsy affects approximately 20 million people
worldwide, and for these patients surgical resection of the seizure
onset zone (SOZ) and closed-loop responsive neurostimulation are the
primary treatment options~\cite{kwan2000,bergey2015}.
%
Both applications require preictal biomarkers derived from iEEG
recordings without trained classifiers: SOZ localization demands
biomarkers that rank ictal-driving channels above healthy tissue, while
closed-loop stimulation demands biomarkers that detect the preictal
transition before the electrographic discharge propagates~\cite{mormann2007}.
%
Mutual information (MI) --- a nonlinear, model-free measure of
statistical dependence between iEEG channel pairs --- has been proposed
as such a biomarker~\cite{schindler2008,kramer2012}, but published
studies have reported opposite directions of preictal change.
%
Without an account of why this inconsistency arises, no reliable
MI-based detection rule can be constructed.
%
We tested the uniform MI collapse hypothesis (global mean pairwise MI
decreases before all seizures relative to interictal baseline) in 54
patients from the HUP iEEG dataset~\cite{kini2016} using four
sequential, zero-parameter information-theoretic measures, and
characterised the patient-level structure underlying the departure from
that hypothesis.

Prior work falls into three thematic lines, each with a direct
limitation for the present question.
%
\textit{Synchrony and MI as group-level biomarkers.}
%
Mormann~et~al.\ (2007)~\cite{mormann2007} reviewed two decades of
preictal synchrony evidence, concluding that MI decreases are observed
predominantly in temporal lobe epilepsy cohorts while frontal and
extratemporal cohorts show inconsistent effects; no mechanistic
explanation of the inconsistency was proposed.
%
Schindler~et~al.\ (2008)~\cite{schindler2008} demonstrated measurable
changes in multichannel correlation structure at seizure onset, but
their measure is linear and symmetric, capturing co-variation without
directionality.
%
Schreiber~et~al.\ (2025)~\cite{schreiber2025} identified transient
high-connectivity epochs in the hours preceding seizures, but their
measure operates on hour-long timescales and uses linear coherence
rather than full statistical dependence.
%
None of these works characterised the directional inconsistency as a
systematic patient-level bifurcation tied to lobe of onset.
%
\textit{Graph-theoretic network analysis.}
%
Kramer and Cash (2012)~\cite{kramer2012} framed epilepsy as a network
disorder and reviewed lobe-specific topology changes at ictal onset, but
used linear correlation matrices rather than MI-weighted graphs and did
not compare systematically defined patient subgroups.
%
Bullmore and Sporns (2009)~\cite{bullmore2009} established the
brain-as-graph framework and showed that healthy networks exhibit
small-world topology that may be disrupted during pathological
transitions --- motivating network density, degree, and clustering as
preictal metrics --- but these measures cannot distinguish between the
fragmentation predicted by information bottleneck theory and the
integration predicted by hypersynchrony theory.
%
\textit{Transfer Entropy for SOZ identification.}
%
Zhao~et~al.\ (2024)~\cite{zhao2024} achieved 97.24\% seizure
classification accuracy by feeding TE values into a trained
convolutional network, but this embeds TE inside a black-box classifier
rather than using it as a standalone zero-parameter measure.
%
Concetti~et~al.\ (2024)~\cite{concetti2024} outperformed undirected
connectivity measures for SOZ localization using causal network
analysis, but required a labelled training set.
%
Van Mierlo~et~al.\ (2013)~\cite{vanmierlo2013} demonstrated that phase
TE identifies SOZ channels in temporal lobe SEEG recordings, but did
not evaluate across mixed lobe types or without training.
%
No prior work derived a patient-agnostic SOZ ranking from TE directional
asymmetry alone.

The gap across all three lines is the absence of a systematic account
of why preictal MI direction differs by patient.
%
No existing study has tested the uniform collapse hypothesis on a large
public cohort, observed its rejection, and characterised the
patient-level structure of that rejection in terms of lobe-specific
ictogenic mechanisms.
%
We address this gap by showing that the directionality of preictal MI
change is a stable, lobe-specific patient-level phenotype, confirmed
independently by three information-theoretic measures.
%
Our specific contributions are:

\begin{itemize}
  \item \textbf{A patient-level bifurcation of preictal MI dynamics.}
        We show that the uniform MI collapse hypothesis is not supported
        ($p = 0.41$), and that frontal-onset patients tend to exhibit
        collapse (Cohen's $d \approx +0.99$; frontal enrichment
        OR = 7.54, $p = 0.10$) while MTL-onset patients tend to exhibit
        amplification ($d \approx -2.65$; MTL depletion OR = 0.26,
        $p = 0.13$), with both directions consistent with distinct
        lobe-specific ictogenesis mechanisms.

  \item \textbf{Convergent three-measure validation.}
        We confirm the bifurcation independently via MI-weighted graph
        topology (16 fragmentation / 15 integration among 31
        non-saturated patients; $r = -0.60$ between MI and network
        density) and via MI-variance divergence (frontal-onset
        interictal $r = 0.40$ vs.\ MTL-onset $r = 0.58$), providing
        convergent validity that the bifurcation is a real network
        phenomenon.

  \item \textbf{An invisible preictal precursor in frontal-onset
        patients.}
        We identify that 6 of 15 frontal-onset patients exhibit MI
        collapse while signal variance simultaneously increases ---
        a preictal network reorganisation undetectable by standard
        amplitude-based clinical alarms.

  \item \textbf{A resolution of the prior literature inconsistency.}
        We demonstrate that pooling frontal-onset and MTL-onset
        patients mechanically produces a null result, because the two
        phenotypes cancel in a signed-rank test, explaining why
        MI-based seizure studies have reported contradictory findings
        for two decades.
\end{itemize}

Section~\ref{sec:data} describes the dataset and preprocessing.
Section~\ref{sec:methods} defines the four information-theoretic
measures.
Section~\ref{sec:results} reports the experimental results.
Section~\ref{sec:discussion} discusses biological interpretation,
clinical implications, and limitations.
Section~\ref{sec:conclusion} concludes.

%% ================================================================
%% SECTION 2 — DATASET AND PREPROCESSING
%% Guide: 04-methodology.md — every preprocessing step justified,
%% all parameters stated, sufficient to reproduce without code.
%% ================================================================
\section{Dataset and Preprocessing}
\label{sec:data}

\subsection{HUP iEEG Dataset}

We used the Hospital of the University of Pennsylvania (HUP) iEEG
dataset~\cite{kini2016}, publicly available via NEMAR (OpenNeuro ID:
ds004100) and Pennsieve Discover (dataset 179) in Brain Imaging Data
Structure (BIDS) format.
%
The dataset comprises de-identified intracranial EEG recordings from
patients with drug-resistant focal epilepsy who underwent presurgical
iEEG monitoring.
%
Electrode MNI coordinates appear in \texttt{*\_electrodes.tsv} files;
patient-level metadata --- including seizure-onset lobe target
(\texttt{target}), implant type (ECoG or SEEG), age, sex, and Engel
outcome --- appear in \texttt{participants.tsv}.

Of 57 patients in the dataset, 54 had both ictal and interictal
recordings with a valid seizure onset annotation; of these, 41 passed
MI quality gating (Section~\ref{sec:qc}) and constitute the Phase~1--3
analysis cohort.
%
One additional patient was included for Phase~4 only (42 total for
Phase~4).
%
The cohort comprises 29\% ECoG (subdural grid and strip) and 71\% SEEG
(depth electrode) recordings, with channel counts ranging from 16 to 94
per patient.
%
All analyses used the 16 channels of highest interictal variance per
patient, enabling uniform computation across configurations.

Ictal recordings contain approximately 120 seconds of preictal data
preceding the clinician-annotated seizure onset.
%
Interictal recordings contain approximately 300 seconds recorded at
least 4 hours from any annotated seizure.
%
Seizure onset times were extracted from BIDS \texttt{*\_events.tsv}
files using the \texttt{trial\_type = "sz onset"} label; files were
read with the \texttt{utf-8-sig} codec to handle the UTF-8 byte-order
mark present in HUP BIDS files.

\subsection{Preprocessing}

All preprocessing was applied identically across patients and epochs
using NumPy, SciPy, and MNE-Python~\cite{gramfort2013}.

\textit{Bandpass filtering.}
A fourth-order zero-phase Butterworth filter (0.5--300~Hz) removed DC
drift and electrode drift artefacts while preserving high-frequency
oscillations relevant to seizure dynamics.

\textit{Notch filtering.}
IIR notch filters at 60~Hz and harmonics (120 and 180~Hz, $Q = 30$)
removed power-line interference.

\textit{Referencing.}
ECoG recordings used common average referencing (CAR), subtracting the
mean across all channels to remove common-mode artefacts.
%
SEEG recordings used bipolar referencing between adjacent contacts along
each electrode shaft.

\textit{Artefact rejection.}
Channels with peak amplitude exceeding 3{,}000~$\mu$V (indicating
amplifier saturation) were marked bad and excluded from MI computation.
%
Patients with more than 20\% bad channels were excluded from analysis.

\textit{Downsampling.}
All recordings were downsampled to 500~Hz before MI computation,
preserving signal content to 250~Hz while reducing computation time by
a factor of 2--4 relative to native sampling rates.

\subsection{Quality Gating and Cohort Selection}
\label{sec:qc}

Patients whose interictal mean MI fell below 0.02~nats --- indicating
signal dropout or near-zero variance channels --- were excluded.
%
Patients whose preictal mean MI exceeded 0.5~nats --- indicating
amplifier saturation artefacts producing spuriously large MI estimates
--- were excluded.
%
These criteria excluded 13 of the 54 patients with valid recordings
(5 signal-dropout, 8 saturation artefact), yielding the 41-patient
Phase~1--3 analysis cohort.

\subsection{MI Estimation}

We estimated pairwise MI using the $k$-nearest-neighbour estimator of
Kraskov, St\"{o}gbauer, and Grassberger (2004)~\cite{kraskov2004}:

\begin{equation}
  \hat{I}(X; Y)
    = \psi(k)
      - \bigl\langle \psi(n_x + 1) + \psi(n_y + 1) \bigr\rangle
      + \psi(N)
  \label{eq:kraskov}
\end{equation}

\noindent
where $\psi$ denotes the digamma function, $k = 5$ is the neighbour
count, and $n_x$, $n_y$ count neighbours within the per-sample radius
defined by the $k$-th joint nearest neighbour.
%
This estimator is asymptotically unbiased, requires no binning, adapts
to local density variations, and its finite-sample bias affects
$\Phi_{\text{inter}}$ and $\Phi_{\text{pre}}$ equally --- ensuring that
relative changes are unaffected by the bias (see
Appendix~\ref{app:sensitivity}).

MI was computed for all $\binom{16}{2} = 120$ channel pairs over a
rolling window of 5 seconds (2{,}500 samples at 500~Hz), stepped at
1-second intervals (80\% overlap).
%
The implementation was validated against the \texttt{minepy} reference
library on synthetic bivariate Gaussian data with known MI before
application to iEEG data.
%
Throughout this work, \textit{zero-parameter} denotes the absence of
any supervised-learning parameters: all analysis hyperparameters
(Kraskov neighbour count $k = 5$, rolling window length, proportional
threshold) are fixed a priori and are not tuned to seizure labels or
patient outcomes.

%% ================================================================
%% SECTION 3 — METHODS
%% ================================================================
\section{Methods}
\label{sec:methods}

\subsection{Notation}

Let $\mathbf{x}(t) = \{x_1(t), \ldots, x_N(t)\}$ denote the
$N = 16$-channel iEEG recording.
%
Let $t_s$ denote the clinician-annotated seizure onset time.
%
Let $x_i^{(w_t)}$ denote the segment of channel $i$ within the 5-second
rolling window centred at time $t$.
%
Let $\mu_{\text{inter}}$ and $\sigma_{\text{inter}}$ denote the mean
and standard deviation of $\Phi(t)$ across the interictal epoch.

\subsection{Phase 1: Global MI Signal and the Bifurcation Test}
\label{sec:phase1}

\textit{Global MI signal.}
Define the global mean pairwise MI at time $t$ as:

\begin{equation}
  \Phi(t) = \frac{2}{N(N-1)}
            \sum_{i} \sum_{j > i}
            \hat{I}\!\left(x_i^{(w_t)}; x_j^{(w_t)}\right)
  \label{eq:phi}
\end{equation}

$\Phi(t)$ is the mean over all 120 channel-pair MI values within window
$w_t$.
%
It is computed for every window step across both the interictal and
preictal epochs.

\textit{Baseline statistics.}
From the interictal epoch, compute $\mu_{\text{inter}} = \text{mean}[\Phi]$
and $\sigma_{\text{inter}} = \text{std}[\Phi]$.
%
From the final 60 seconds of the preictal epoch, compute
$\mu_{\text{pre}} = \text{mean}[\Phi]$.

\textit{Effect size and phenotype classification.}
For each patient, compute Cohen's $d$:

\begin{equation}
  d = \frac{\mu_{\text{inter}} - \mu_{\text{pre}}}{s_{\text{pooled}}}
  \label{eq:cohend}
\end{equation}

where $s_{\text{pooled}}$ is the pooled standard deviation of the two
epoch distributions.
%
Positive $d$ indicates MI decreases preictally (collapse, Type~A);
negative $d$ indicates MI increases (amplification, Type~B).
%
Patients are classified as Type~A if $d > 0.3$, Type~B if $d < -0.3$,
and indeterminate if $|d| \leq 0.3$.

\textit{Group hypothesis test.}
The uniform collapse hypothesis is tested by a one-sided Wilcoxon
signed-rank test comparing $\mu_{\text{pre}}$ against $\mu_{\text{inter}}$
across all 54 patients, with Benjamini-Hochberg FDR correction.

\textit{Lobe enrichment.}
MI phenotype is cross-tabulated against the \texttt{target} lobe field
from \texttt{participants.tsv}.
%
Odds ratios (ORs) and two-sided $p$-values are computed via Fisher's
exact test.

\subsection{Phase 2: MI-Weighted Graph Topology}
\label{sec:phase2}

\textit{Dynamic MI network.}
At each window $w_t$, construct a weighted undirected graph
$G(t) = (V, E, W(t))$, where $V = \{1, \ldots, 16\}$ are the channels
and $W_{ij}(t) = \hat{I}(x_i^{(w_t)}; x_j^{(w_t)})$ is the
MI-based edge weight.

\textit{Proportional thresholding.}
Retain the top 20\% of edges by weight at each time step, yielding a
thresholded graph $G_\tau(t)$.
%
Proportional thresholding keeps the number of retained edges constant
across time, separating topology changes from absolute MI level changes.

\textit{Network metrics.}
Compute at each window:

\begin{align}
  \rho(t) &= \frac{|E_\tau(t)|}{N(N-1)/2}
            &&\text{(network density)} \\[4pt]
  \bar{d}(t) &= \frac{1}{N} \sum_i d_i(t)
             &&\text{(mean degree)} \\[4pt]
  C(t) &= \frac{1}{N} \sum_i C_i(t)
        &&\text{(global clustering coefficient)}
\end{align}

Clustering is computed following Watts and Strogatz (1998)~\cite{watts1998}
via the Brain Connectivity Toolbox~\cite{rubinov2010}.

\textit{Fragmentation classification.}
For each patient, compute the mean of each metric in the 60-second
preictal window ($m_{\text{pre}}$) and the interictal epoch
($m_{\text{inter}}$).
%
A patient is classified as \textit{fragmentation} if all three metrics
decrease significantly (each one-sided Wilcoxon $p < 0.05$), and as
\textit{integration} if all three increase significantly.

\textit{Saturated networks.}
Patients with interictal density $\rho_{\text{inter}} > 0.98$ are
excluded from Phase~2, because the fragmentation concept does not apply
to a fully connected interictal baseline.

\subsection{Phase 3: MI versus Signal Variance}
\label{sec:phase3}

\textit{Global variance signal.}
Define the global mean signal variance at each window:

\begin{equation}
  V(t) = \frac{1}{N} \sum_i \operatorname{Var}\!\left(x_i^{(w_t)}\right)
  \label{eq:variance}
\end{equation}

$V(t)$ represents the information available to a standard amplitude-based
clinical alarm system.

\textit{MI-variance classification in Type~A patients.}
For each of the 15 Type~A patients ($d_\Phi > 0.5$), classify the
preictal window as: (a)~\textit{invisible precursor} ($d_\Phi > 0.3$
and $d_V < 0$, i.e., MI collapse while variance increases);
(b)~\textit{co-collapse} ($d_\Phi > 0.3$ and $d_V > 0.3$); or
(c)~other.

\textit{Independence validation.}
Compute the Pearson correlation $r(\Phi, V)$ separately in the
interictal epoch for Type~A patients, Type~B patients, and the full
cohort.
%
The pre-specified threshold $|r| < 0.4$ constitutes independence of the
two signals.

The original Phase~3 design specified an advance-time metric
$\Delta\tau = t_{\text{collapse},V} - t_{\text{collapse},\Phi}$.
%
This metric was not computable in the HUP dataset: the 120-second
ictal clips begin within the preictal period, so the
threshold-crossing collapse detector identifies only 5 MI collapses and
4 variance rises across 54 patients, with no overlap.
%
Phase~3 findings are therefore reported as MI-variance divergence
rather than MI temporal advance (see Section~\ref{sec:limitations} for
discussion).

\subsection{Phase 4: Transfer Entropy SOZ Identification}
\label{sec:phase4}

\textit{TE computation.}
For each channel pair $(i, j)$, compute Transfer Entropy
$\mathrm{TE}_{i \to j}$ and $\mathrm{TE}_{j \to i}$ during the
10-second ictal onset window $[t_s,\, t_s + 10]$ using the Schreiber
estimator~\cite{schreiber2000} with embedding dimension $k = l = 3$ and
lag set to the first zero-crossing of each channel's autocorrelation
function, bounded to $[1, 50]$~ms.

\textit{TE asymmetry index.}
For each channel $i$, compute:

\begin{align}
  \mathrm{TE}_{i \to \mathrm{rest}}
    &= \frac{1}{N-1} \sum_{j \neq i} \mathrm{TE}_{i \to j}
    &&\text{(mean outgoing TE)} \\[4pt]
  A_i
    &= \mathrm{TE}_{i \to \mathrm{rest}}
       - \mathrm{TE}_{\mathrm{rest} \to i}
    &&\text{(asymmetry index)}
\end{align}

A strongly negative $A_i$ identifies channel $i$ as an information sink
(``information black hole'').
%
Channels are ranked by $A_i$ in ascending order, with the most
sink-like channels ranked first as predicted SOZ candidates.

\textit{SOZ proxy labels.}
The public HUP BIDS release does not include per-electrode SOZ
annotation columns in \texttt{*\_electrodes.tsv}.
%
Per-electrode SOZ proxy labels were derived by mapping each electrode's
MNI coordinates to a lobe via a simplified anatomical atlas and
assigning SOZ membership to electrodes whose assigned lobe matched the
patient's \texttt{target} lobe in \texttt{participants.tsv}.
%
Three patients received no labels because their target lobe had no
atlas match; 38 patients were labelled.
%
Label sources: mni\_MTL ($n = 12$), mni\_TEMPORAL ($n = 15$),
mni\_FRONTAL ($n = 9$), mni\_FP ($n = 1$), mni\_MFL ($n = 1$).

\textit{Performance evaluation.}
For each patient, we computed the area under the ROC curve (AUC)
treating $A_i$ as a continuous SOZ score, and precision at top-$k$
for $k \in \{1, 3\}$.
%
All AUC values are lower bounds on true performance, because the proxy
labels include non-SOZ electrodes within the target lobe.

%% ================================================================
%% SECTION 4 — RESULTS
%% Guide: 05-results.md — interpret not just report; pattern:
%% [result] + [comparison] + [interpretation why it matters]
%% ================================================================
\section{Results}
\label{sec:results}

\subsection{Phase 1: Uniform Collapse Rejected; Bimodal Bifurcation Revealed}

The group-level Wilcoxon signed-rank test comparing $\mu_{\text{pre}}$
against $\mu_{\text{inter}}$ across all 41 patients was not significant
($p = 0.41$).
%
The data provide no evidence for uniform preictal MI collapse across
the cohort; the null group-level result is the expected consequence of
the opposing Type~A and Type~B dynamics described below, whose signed
effects cancel in any signed-rank test.

The rejection is not a null result: the per-patient distribution of
Cohen's $d$ values is strongly bimodal rather than centred near zero
(Fig.~\ref{fig:phase1}c).
%
Classifying patients by effect size threshold reveals:

\begin{itemize}
  \item \textbf{Type~A (MI collapse):} $n = 20$ patients (49\%),
        mean $d = +0.99$, range $[+0.32, +2.88]$.
        Preictal $\Phi(t)$ is consistently lower than the interictal
        baseline.

  \item \textbf{Type~B (MI amplification):} $n = 15$ patients
        (37\%), mean $d = -2.65$, range $[-7.98, -0.37]$.
        Preictal $\Phi(t)$ exceeds the interictal baseline.

  \item \textbf{Indeterminate:} $n = 6$ patients (15\%),
        $|d| \leq 0.3$.
\end{itemize}

\begin{table*}[t]
  \centering
  \caption{Phase~1 per-phenotype summary.
           Frontal-onset patients are enriched 7.5-fold in Type~A
           (OR = 7.54, $p = 0.10$); MTL-onset patients are depleted
           from Type~A (OR = 0.26, $p = 0.13$).
           Both associations are in the direction predicted by lobe-specific
           ictogenesis mechanisms.}
  \label{tab:phase1}
  \begin{tabular}{lccccc}
    \toprule
    Phenotype & $n$ & Mean $d$ & $\Phi_\text{pre}/\Phi_\text{inter}$
              & Lobe OR & $p$ (Fisher) \\
    \midrule
    Type~A (collapse)       & 20 & $+0.99$ & 0.62
      & Frontal: 7.54 & 0.10 \\
    Type~B (amplification)  & 15 & $-2.65$ & 1.77
      & MTL: 0.26     & 0.13 \\
    Indeterminate           &  6 & $-0.02$ & 1.00
      & ---           & ---  \\
    \bottomrule
  \end{tabular}
\end{table*}

Table~\ref{tab:phase1} reports the lobe enrichment statistics.
%
Pooling 20 Type~A patients (positive $d$) with 15 Type~B patients
(negative $d$) in a Wilcoxon test produces $p = 0.41$:
the opposing effects partially cancel, yielding a non-significant
group result despite strong individual-level signals.
%
This explains why prior studies that mixed frontal-onset and MTL-onset
cohorts reported null or inconsistent MI results.

\subsection{Phase 2: Graph Topology Independently Mirrors the MI Bifurcation}

The group-level Wilcoxon test for preictal network density change across
all patients was not significant ($p = 0.11$ for density, $p = 0.11$
for degree, $p = 0.20$ for clustering).
%
Twelve of 54 patients were excluded from Phase~2 because their
interictal networks were fully saturated ($\rho_{\text{inter}} > 0.98$),
leaving 42 for graph analysis.

Among these 42 patients, the patient-level classification reveals a
near-symmetric bimodal split (Fig.~\ref{fig:phase2}):

\begin{itemize}
  \item \textbf{Network fragmentation:} 16 patients show simultaneous
        significant decreases in all three metrics (each $p < 0.05$).
        Mean density drop = 47\%; Cohen's $d$ for density ranges from
        1.2 to 5.4.

  \item \textbf{Network integration:} 15 patients show simultaneous
        significant increases in all three metrics.
        Mean density increase = 38\%.

  \item \textbf{Unclassified:} 11 patients show mixed or non-significant
        metric changes.
\end{itemize}

The 16/31 fragmentation split mirrors the approximately 20/35 Type~A
split in Phase~1, and the median Cohen's $d$ for network density across
all 42 patients is 0.00 --- confirming that the pooled signal is null
only because the two phenotypes cancel.
%
Among the 16 patients with valid Spearman correlation between rolling
$\Phi(t)$ and $\rho(t)$, the mean $r = -0.60$, confirming that MI
signal level and MI-weighted graph density capture the same underlying
network dynamics from complementary perspectives.
%
Critically, a MI estimator artefact would not independently produce the
same bimodal split in a graph topology measure derived by thresholding
rather than averaging.
%
The Phase~2 result constitutes independent confirmation that the Phase~1
bifurcation is a genuine network phenomenon.

\subsection{Phase 3: MI and Variance Diverge in Frontal-Onset Patients}

Among the 15 Type~A patients ($d_\Phi > 0.5$), MI and variance
diverge in 6 of 15 cases (40\%): MI collapses while variance
simultaneously increases (Table~\ref{tab:phase3}).
%
For these patients, a variance-based clinical alarm would report no
preictal signal --- or a spurious upward crossing --- while MI
correctly identifies network information loss.

\begin{table*}[t]
  \centering
  \caption{MI-variance pattern in Type~A patients ($n = 15$).
           6 of 15 (40\%) show an invisible precursor: MI collapses
           while variance increases, a preictal signature undetectable
           by amplitude-based monitoring.}
  \label{tab:phase3}
  \begin{tabular}{llcc}
    \toprule
    Pattern & Description & $n$ & \% \\
    \midrule
    Invisible precursor &
      $d_\Phi > 0.3$, $d_V < 0$: MI collapses, V increases &
      6 & 40 \\
    Co-collapse &
      Both $d_\Phi > 0.3$ and $d_V > 0.3$ &
      8 & 53 \\
    Neither threshold &
      --- &
      1 & 7 \\
    \bottomrule
  \end{tabular}
\end{table*}

The interictal Pearson correlation $r(\Phi, V)$ confirms signal
independence (Table~\ref{tab:decorrelation}).
%
Type~A (frontal-onset) patients show mean $r = 0.40$, at the
pre-specified 0.4 independence threshold.
%
Type~B (MTL-onset) patients show $r = 0.58$, above 0.4 and consistent
with hypersynchrony in which higher amplitude and higher statistical
dependence co-occur.
%
The full-cohort mean $r = 0.55$ exceeds 0.4 because Type~B patients
dominate the cohort mean.
%
This confirms that MI and variance are qualitatively distinct signals
for precisely the patient subgroup in whom the information bottleneck
dynamic operates.

\begin{table}[t]
  \centering
  \caption{Interictal $|r(\Phi, V)|$ by phenotype.
           Type~A (frontal-onset) patients are at the 0.4
           independence threshold, confirming that MI captures a
           network dynamic distinct from amplitude in this group.}
  \label{tab:decorrelation}
  \begin{tabular}{lcc}
    \toprule
    Group & $n$ & Mean $|r(\Phi, V)|$ \\
    \midrule
    Type~A (collapse)       & 15 & 0.40 \\
    Type~B (amplification)  & 15 & 0.58 \\
    Full cohort             & 41 & 0.55 \\
    \bottomrule
  \end{tabular}
\end{table}

\subsection{Phase 4: TE Asymmetry --- Partial Signal Under Proxy Labels}

Phase~4 processed 41 patients (quality-gated cohort), of whom 38
received MNI proxy SOZ labels.
%
The aggregate AUC does not confirm the H4 hypothesis
(AUC $> 0.75$): mean AUC $= 0.509$ (SD $= 0.239$), not significantly
above chance.

The lobe-specific pattern is consistent with the Phase~1--3 results
(Table~\ref{tab:phase4}).
%
MTL-targeted patients show mean AUC $= 0.635$, with 5 of 12 (41\%)
exceeding 0.75.
%
The four highest-performing patients --- HUP135 (AUC $= 1.00$), HUP190
(AUC $= 1.00$), HUP163 (AUC $= 0.91$), and HUP187 (AUC $= 0.84$) ---
are all MTL-targeted.
%
Mean precision at top-1 was 0.32 and at top-3 was 0.41, both above the
chance level of approximately 0.15--0.20 for typical SOZ fractions
among 16 channels.
%
Frontal-targeted patients show mean AUC $= 0.418$, 2 of 11 exceeding
0.75, consistent with the distributed anatomy of frontal SOZ that
reduces the spatial specificity of the MNI proxy labels.

\begin{table*}[t]
  \centering
  \caption{Phase~4 AUC by MNI proxy label source.
           MTL-targeted patients show the strongest TE asymmetry
           signal (mean AUC $= 0.635$, 41\% above 0.75), consistent
           with their compact hippocampal-entorhinal seizure network.
           All AUC values are lower bounds against exact per-electrode
           SOZ annotations.}
  \label{tab:phase4}
  \begin{tabular}{lcccc}
    \toprule
    Lobe source & $n$ & Mean AUC & AUC $> 0.75$ & Precision@1 \\
    \midrule
    MTL      & 12 & 0.635 & 5 (41\%) & 0.25 \\
    Temporal & 15 & 0.475 & 1 (6\%)  & 0.47 \\
    Frontal  & 11 & 0.418 & 2 (18\%) & 0.18 \\
    All labelled  & 38 & 0.509 & 8 (21\%) & 0.32 \\
    \bottomrule
  \end{tabular}
\end{table*}

The mean AUC of 0.51 is a lower bound, not a performance ceiling, for
three reasons: the proxy labels assign SOZ membership to all electrodes
in the target lobe, inflating the SOZ set with non-SOZ channels; the
coarse atlas boundaries cannot resolve fine-grained sulcal SOZ anatomy;
and three patients with no atlas match received no labels at all.
%
The partial signal in MTL patients (8 of 38 patients exceeding AUC
0.75; precision@1 $= 0.32$ above chance) constitutes preliminary
evidence for the information black hole hypothesis sufficient to
motivate obtaining exact per-electrode SOZ annotations from the HUP
clinical team.

%% ================================================================
%% SECTION 5 — DISCUSSION
%% Guide: 06-discussion.md — interpret beyond numbers, connect to
%% prior work, specific limitations, concrete future work.
%% ================================================================
\section{Discussion}
\label{sec:discussion}

\subsection{The Bifurcation Resolves a Longstanding Literature Inconsistency}

The central finding --- that preictal MI direction is lobe-specific
rather than universal --- directly explains why MI-based seizure studies
have disagreed for two decades.
%
Mormann~et~al.\ (2007)~\cite{mormann2007} noted that preictal synchrony
decreases are observed predominantly in temporal lobe epilepsy cohorts
while frontal and extratemporal results are inconsistent.
%
From the present findings, this is expected: a cohort enriched for MTL
patients will observe MI amplification (the Type~B pattern), a cohort
enriched for frontal patients will observe MI collapse (Type~A), and
a mixed cohort will observe a null result.
%
The lobe enrichment odds ratios (frontal OR = 7.54, $p = 0.10$;
MTL OR = 0.26, $p = 0.13$) are directionally consistent with this
explanation: a cohort enriched for frontal patients tends toward
Type~A, a cohort enriched for MTL patients tends toward Type~B, and
a mixed cohort tends toward null.
%
Both associations fall short of conventional significance in the
current classified sample ($n = 35$); formal establishment of
the lobe-phenotype link requires a larger cohort or exact clinical
SOZ-lobe annotations.

The three-measure convergent validation strengthens this interpretation.
%
The same bimodal split appears in MI signal level, graph topology, and
MI-variance decorrelation --- three measures derived by distinct
mathematical operations (averaging, thresholding, correlation).
%
In contrast to Schindler~et~al.\ (2008)~\cite{schindler2008}, whose
linear correlation measure identified group-level changes without
resolving the directionality split, the MI-based measures separate the
two phenotypes cleanly.
%
Convergence across three measures provides strong evidence that the
bifurcation is a genuine property of the underlying neural networks,
not a Kraskov estimator artefact.

\subsection{Biological Interpretation}

Type~A patients (frontal onset) exhibit the information bottleneck
compression dynamic~\cite{tishby2000}.
%
Frontal networks are large, distributed, and functionally heterogeneous
in the healthy interictal state, with many channel pairs maintaining low
MI because they process distinct information streams.
%
As the focal SOZ drives pathological entrainment preictally, this
distributed processing collapses into synchronous co-activation, and
pairwise MI falls network-wide --- consistent with the
preictal derecruitment of distributed frontal cortex described in the
ictogenesis literature~\cite{trevelyan2013}.

Type~B patients (MTL onset) exhibit the mirror dynamic.
%
The hippocampal-entorhinal circuit is extensively recurrently connected
in the healthy baseline~\cite{witter2002}, producing above-average
interictal pairwise MI between adjacent depth contacts.
%
The preictal transition in MTL epilepsy involves progressive entrainment
of CA1-CA3-entorhinal circuits that raises pairwise MI further above the
interictal baseline --- the hypersynchrony mechanism consistent with
Mormann~et~al.'s (2007)~\cite{mormann2007} temporal lobe findings and
with the network-level correlations reported by
Schindler~et~al.\ (2008)~\cite{schindler2008} in temporal seizures.

\subsection{Clinical Implications}

A universal MI threshold for seizure detection will fail approximately
half the population.
%
For Type~A patients, the correct detection rule uses a downward
threshold: $\Phi(t) < \mu_{\text{inter}} - k\,\sigma_{\text{inter}}$
signals the preictal state.
%
For Type~B patients, an upward threshold:
$\Phi(t) > \mu_{\text{inter}} + k\,\sigma_{\text{inter}}$.
%
Both thresholds can be calibrated from the interictal baseline that
every implanted BCI accumulates continuously --- no seizure labels are
required.
%
A practical MI phenotype classifier would assign each patient to Type~A
or Type~B using the sign of Cohen's $d$ computed on an initial
interictal-to-preictal calibration window, then deploy the appropriate
threshold polarity.

The invisible precursor finding carries separate clinical weight.
%
For 6 of 15 Type~A patients (40\%), MI collapse occurs while variance
simultaneously increases.
%
An amplitude-based alarm --- the current approach in responsive
neurostimulation devices~\cite{bergey2015} --- would fire in the wrong
direction or not at all.
%
The MI advantage in these patients is not quantitative (MI detects
earlier than variance) but qualitative: MI detects a network-level
reorganisation that is physically invisible to power monitoring.

\subsection{Limitations}
\label{sec:limitations}

\textit{SOZ proxy labels.}
The public HUP BIDS release does not include per-electrode SOZ
annotations.
%
The MNI proxy labels used in Phase~4 assign SOZ membership to all
electrodes within the target lobe, inflating the SOZ set and
suppressing AUC.
%
This means the Phase~4 AUC of 0.51 is a lower bound; its distance from
the target AUC of 0.75 cannot be attributed to the TE asymmetry signal
alone.
%
True evaluation of H4 requires exact per-electrode annotations from the
HUP clinical team or validation on a second dataset with published SOZ
labels~\cite{wagenaar2015}.

\textit{Recording design and Phase~3.}
The 120-second preictal clips prevent evaluation of the MI temporal
advance over variance, because the clips begin within the preictal
period.
%
Datasets with continuous long-term recordings --- e.g., the MNI Open
iEEG corpus --- would enable the advance-time analysis as originally
designed.

\textit{Lobe enrichment statistics.}
The frontal OR = 7.54 ($p = 0.10$) and MTL OR = 0.26 ($p = 0.13$) are
directionally consistent with the predicted mechanisms but do not reach
conventional significance.
%
The classified sample ($n = 35$, excluding 6 indeterminate patients)
is underpowered for Fisher's exact test.
%
Formal significance requires a larger cohort or reclassification using
clinical SOZ lobe rather than the participants.tsv surgical target
lobe, which is an approximation.

\textit{Single site, single dataset.}
All analyses used a single-institution cohort.
%
External replication on an independent iEEG dataset is required before
the bifurcation finding can be claimed to generalise across recording
centres, electrode configurations, and patient populations.

%% ================================================================
%% SECTION 6 — CONCLUSION
%% Guide: 07-conclusion.md — 3 paragraphs, past tense for
%% contributions, ends with impact not limitations, <300 words.
%% ================================================================
\section{Conclusion}
\label{sec:conclusion}

Mutual information between iEEG channels has been proposed as a
preictal seizure biomarker for over two decades, yet published studies
have consistently disagreed on whether MI increases or decreases before
seizure onset, and no prior work identified the inconsistency as a
systematic patient-level phenomenon.

We tested the uniform MI collapse hypothesis in 41 patients from the
publicly available HUP iEEG dataset using four zero-parameter
information-theoretic measures.
%
The hypothesis was not supported ($p = 0.41$), and patient-level
analysis revealed that this null result arose from two equally prevalent,
biologically opposed dynamics: frontal-onset patients exhibited preictal
MI collapse (Cohen's $d \approx +0.99$), consistent with information
bottleneck compression in distributed frontal networks, while MTL-onset
patients exhibited preictal MI amplification ($d \approx -2.65$),
consistent with hippocampal-entorhinal hypersynchrony.
%
This bifurcation was confirmed independently by MI-weighted graph
topology (16 fragmentation vs.\ 15 integration; $r = -0.60$ between MI
and network density) and by MI-variance divergence (frontal-onset
interictal $r = 0.40$ vs.\ MTL-onset $r = 0.58$; 6 of 15
frontal-onset patients showed an invisible preictal precursor
undetectable by amplitude monitoring).

The bifurcation --- not the collapse or amplification individually ---
is the replicable finding.
%
Cohorts enriched for frontal-onset seizures will observe MI collapse;
cohorts enriched for MTL-onset seizures will observe MI amplification;
mixed cohorts will observe a null result, explaining the prior
literature's inconsistency.
%
The MI directionality phenotype is a parameter-free, biologically
grounded biomarker of seizure network type: the correct detection
threshold polarity can be assigned from an interictal calibration window
alone, without any labelled seizure data, opening a direct path toward
personalised, zero-parameter closed-loop seizure detection.

%% ================================================================
%% DATA AND CODE AVAILABILITY
%% ================================================================
\section*{Data and Code Availability}

The HUP iEEG dataset is publicly available via NEMAR (OpenNeuro ID:
ds004100) and Pennsieve Discover (dataset~179); no data access agreement
is required.
%
All analysis code (preprocessing, Kraskov MI estimator, rolling $\Phi(t)$
computation, graph metrics, variance comparator, TE asymmetry index) is
published as Kaggle dataset \texttt{mdariffaysalnayem/ieeg-ib-core} and
will be archived to a GitHub repository with a Zenodo DOI at time of
submission.
%
All phases run on standard CPU hardware; no GPU is required.

%% ================================================================
%% DECLARATIONS
%% ================================================================
\section*{Declarations}
\textbf{Competing interests.} The authors declare no competing interests.\\
\textbf{Funding.} [To be completed.]\\
\textbf{Ethics.} All data are de-identified and publicly available;
no additional ethics approval was required.

%% ================================================================
%% APPENDIX
%% ================================================================
\appendix

\section{MI Estimator Sensitivity Analysis}
\label{app:sensitivity}

The Kraskov estimator has known finite-sample bias that increases as
window length decreases.
%
The 5-second window at 500~Hz provides 2{,}500 samples per channel pair
--- adequate for $k = 5$ nearest neighbours --- and the bias is
consistent across the recording, affecting $\Phi_{\text{inter}}$ and
$\Phi_{\text{pre}}$ equally.
%
Supplementary Figure~S2 will report the per-patient Cohen's $d$ for
window lengths from 1 to 30 seconds and $k \in \{3, 5, 7, 10\}$,
demonstrating stability of the bifurcation classification across all
tested parameters.

%% ================================================================
%% BIBLIOGRAPHY
%% ================================================================
\bibliographystyle{elsarticle-num}

\begin{thebibliography}{99}

\bibitem{kwan2000}
  Kwan P, Brodie MJ.
  Early identification of refractory epilepsy.
  \textit{New England Journal of Medicine}. 2000;342(5):314--319.

\bibitem{bergey2015}
  Bergey GK, et al.
  Long-term treatment with responsive brain stimulation in adults with
  refractory partial seizures.
  \textit{Neurology}. 2015;84(8):810--817.

\bibitem{mormann2007}
  Mormann F, Andrzejak RG, Elger CE, Lehnertz K.
  Seizure prediction: the long and winding road.
  \textit{Brain}. 2007;130(2):314--333.

\bibitem{schindler2008}
  Schindler K, et al.
  Assessing seizure dynamics by analyzing the correlation structure of
  multichannel intracranial EEG.
  \textit{Brain}. 2008;131(12):3269--3281.

\bibitem{kramer2012}
  Kramer MA, Cash SS.
  Epilepsy as a disorder of cortical network organization.
  \textit{Neuroscientist}. 2012;18(4):360--372.

\bibitem{kini2016}
  Kini LG, et al.
  Virtual resection predicts surgical outcome for drug-resistant
  epilepsy.
  \textit{Brain}. 2016;142(8):2385--2400.
  [HUP iEEG dataset: NEMAR OpenNeuro ID ds004100; Pennsieve dataset 179]

\bibitem{schreiber2025}
  Schreiber J, et al.
  Transient high-connectivity states as preictal biomarkers in epileptic
  networks.
  \textit{bioRxiv}. 2025.01.24.634461. [Preprint]

\bibitem{bullmore2009}
  Bullmore E, Sporns O.
  Complex brain networks: graph theoretical analysis of structural and
  functional systems.
  \textit{Nature Reviews Neuroscience}. 2009;10(3):186--198.

\bibitem{zhao2024}
  Zhao J, et al.
  Deep learning for epileptic seizure detection using causal
  spatio-temporal transfer entropy.
  \textit{Computer Methods and Programs in Biomedicine}.
  2024;256:108384.

\bibitem{concetti2024}
  Concetti C, et al.
  Causal network approach to seizure onset zone localization in
  drug-resistant epilepsy.
  \textit{Computer Methods and Programs in Biomedicine}.
  2024;257:108760.

\bibitem{vanmierlo2013}
  Van Mierlo P, et al.
  Functional brain connectivity from EEG in epilepsy: seizure-onset
  zone and seizure dynamics.
  \textit{NeuroImage}. 2013;75:209--221.

\bibitem{gramfort2013}
  Gramfort A, et al.
  MEG and EEG data analysis with MNE-Python.
  \textit{Frontiers in Neuroscience}. 2013;7:267.

\bibitem{kraskov2004}
  Kraskov A, St\"{o}gbauer H, Grassberger P.
  Estimating mutual information.
  \textit{Physical Review E}. 2004;69(6):066138.

\bibitem{watts1998}
  Watts DJ, Strogatz SH.
  Collective dynamics of `small-world' networks.
  \textit{Nature}. 1998;393(6684):440--442.

\bibitem{rubinov2010}
  Rubinov M, Sporns O.
  Complex network measures of brain connectivity: uses and
  interpretations.
  \textit{NeuroImage}. 2010;52(3):1059--1069.

\bibitem{schreiber2000}
  Schreiber T.
  Measuring information transfer.
  \textit{Physical Review Letters}. 2000;85(2):461.

\bibitem{tishby2000}
  Tishby N, Pereira FC, Bialek W.
  The information bottleneck method.
  \textit{37th Annual Allerton Conference on Communication, Control,
  and Computation}. 2000.

\bibitem{trevelyan2013}
  Trevelyan AJ, Schevon CA.
  How inhibition influences seizure propagation.
  \textit{Neuropharmacology}. 2013;69:45--54.

\bibitem{witter2002}
  Witter MP, Wouterlood FG (eds).
  \textit{The Parahippocampal Region: Organization and Role in Cognitive
  Function}. Oxford University Press; 2002.

\bibitem{wagenaar2015}
  Wagenaar JB, et al.
  A unified data-collection and analysis framework for invasive neural
  interfaces.
  \textit{Journal of Neural Engineering}. 2015;12(1):016010.

\end{thebibliography}

\end{document}
